%!TeX root=MemoriaTFG.tex

\chapter{Introducción}
\label{cap:intro}

En el presente Trabajo Final de Grado (TFG) se evalúa de forma
sistemática el rendimiento de los modelos fundacionales disponibles
para imagen dermatológica clínica, sobre un conjunto homogéneo de
protocolos experimentales y datasets externos públicos. Por
\emph{modelo fundacional} se entiende en este trabajo un modelo de
gran escala preentrenado sobre un volumen amplio de datos y
reutilizable, mediante adaptación, en múltiples tareas posteriores;
el estudio evalúa los disponibles con pesos abiertos y, sobre esa
base, construye una adaptación propia.

El trabajo articula así tres aportaciones de naturaleza distinta pero
complementaria: un \emph{benchmark} independiente y reproducible de
catorce modelos fundacionales bajo un protocolo común; un
\emph{dataset} clínico propio en español, DermapixelAI~1.0, entregado
como contribución independiente; y un \emph{modelo adaptado},
Dermapixel R0, que demuestra la viabilidad de especializar un modelo
fundacional preexistente al castellano con recursos modestos. Las tres
se desarrollan en el marco de dos colaboraciones preexistentes con la
Dra.~R.~Taberner del Hospital Universitario Son Llàtzer (HUSLL).

En este capítulo se presenta el contexto clínico del problema, el
marco operativo en el que se inscribe el trabajo, los objetivos
formales que lo articulan y la estructura del resto del documento. El estado del arte específico ---modelos fundacionales
para imagen médica y dermatológica, métodos de adaptación y
evaluación zero-shot--- se desarrolla en el
capítulo~\ref{cap:sota}.


\section{Contexto clínico}
\label{sec:contexto}

La dermatología clínica atiende un espectro amplio de entidades con
distribución concentrada en un número reducido de motivos de
consulta. El estudio observacional DIADERM, sobre $10\,999$
diagnósticos en $8\,953$ pacientes en territorio español,
ordena los diez diagnósticos más prevalentes como queratosis
actínica, carcinoma basocelular, nevus melanocítico, queratosis
seborreica, otras neoplasias benignas, psoriasis, acné, verrugas
víricas, otros trastornos de la pigmentación y otras
dermatitis~\cite{diaderm2018}. El trabajo local de Taberner et~al.\
sobre el servicio dermatológico del HUSLL documenta $213$
diagnósticos diferentes en un único año
asistencial~\cite{taberner2010_motivos}. Ambos coinciden en un peso
elevado de dermatosis benignas, inflamatorias e infecciosas
habituales.

La sensibilidad temporal del proceso diagnóstico resulta crítica en
el melanoma cutáneo. La supervivencia específica a cinco años
desciende desde aproximadamente $99\%$ en estadio~I hasta
$25\%$ en estadio~IV según la octava edición de la American
Joint Committee on Cancer~\cite{ajcc8}. El acceso al especialista
se realiza mediante derivación desde atención primaria. La
capacidad asistencial del circuito presencial es limitada frente a
una demanda creciente, lo que justifica el interés por sistemas de
apoyo al diagnóstico basados en imagen.

El sustrato visual del diagnóstico dermatológico se articula en tres
modalidades: fotografía clínica (presentación general y contexto
anatómico), dermatoscopia (patrones pigmentarios y estructurales no
visibles a simple vista) e histopatología (confirmación
tisular). Sobre esta tríada operan los sistemas de apoyo basados en
aprendizaje automático.


\section{Marco operativo previo}
\label{sec:trabajo-previo}

Las dos colaboraciones preexistentes con la Dra.~R.~Taberner,
dermatóloga del HUSLL, en cuyo marco se desarrolla el trabajo, se
detallan a continuación. La primera es la modernización tecnológica
del archivo Dermapixel,
blog mantenido por la Dra.~Taberner desde 2010. El archivo reúne
más de $700$ casos clínicos comentados con imagen clínica y
dermatoscópica, historia abreviada, diagnóstico y razonamiento
diferencial. En paralelo a este TFG se ha desarrollado su migración
a una pila tecnológica propia. La segunda es un proyecto
institucional del Servei de Salut de les Illes Balears (IB-Salut)
sobre teledermatología hospitalaria. En él se desarrolla una
herramienta de apoyo a la primera valoración dermatológica en el
circuito de derivación desde atención primaria. Ambas líneas
orientan la decisión de priorizar modelos con código y pesos
públicos, por su viabilidad de integración hospitalaria y su
reproducibilidad. Esta prioridad introduce un sesgo de selección
explícito ---deja fuera modelos cerrados potencialmente más capaces,
como DermFM-Zero, cuyos pesos no son públicos
(sección~\ref{sec:lim-pesos})---, asumido a cambio de un marco de
evaluación íntegramente reproducible; los modelos cerrados accesibles
vía API se incorporan únicamente como comparadores contextuales.


\section{Objetivos del trabajo}
\label{sec:objetivos}

Este trabajo evalúa el rendimiento y la capacidad de generalización
de los modelos fundacionales dermatológicos disponibles, en función
del volumen de datos etiquetados, del tipo de supervisión, de la
heterogeneidad taxonómica entre conjuntos de datos y de la disparidad
entre fototipos cutáneos. El resultado de esa evaluación condiciona la
viabilidad de integrarlos en teledermatología hospitalaria con cohortes
mediterráneas en español, escenario donde las referencias recientes del
grupo de la Universidad Monash ---PanDerm, la familia DermLIP sobre
Derm1M y DermFM-Zero--- carecen de evaluación independiente comparable
a la fecha de cierre.

El trabajo se articula en torno a tres objetivos formales,
verificables sobre los resultados experimentales del
capítulo~\ref{cap:resultados}.

\paragraph{O1. Evaluación experimental sistemática de los modelos
fundacionales dermatológicos sobre conjuntos de datos externos.}
Sobre el hardware del proyecto ---una NVIDIA DGX Spark de
arquitectura aarch64, cuya configuración completa se detalla en la
sección~\ref{sec:infraestructura}--- se implementa el conjunto de
modelos fundacionales con código y pesos públicos a la fecha de
inicio. Se evalúa su
comportamiento sobre doce datasets externos
públicos\footnote{HAM10000~\cite{ham10000},
BCN20000~\cite{bcn20000}, PAD-UFES-20~\cite{padufes},
DDI~\cite{ddi}, Dermnet~\cite{dermnet}, Derm7pt
(clínica y dermatoscópica)~\cite{kawahara2019}, HIBA, MSKCC, WSI
patches, Fitzpatrick17k~\cite{fitzpatrick17k},
SkinCon~\cite{skincon} e ISIC2018~\cite{isic2018}.} mediante cuatro
protocolos principales (sondeo lineal, ajuste fino supervisado,
segmentación binaria promptable sobre SAM2.1~\cite{sam2},
clasificación zero-shot multimodal) y tres protocolos
complementarios (comparación con modelos de lenguaje multimodales,
análisis de equidad por fototipo Fitzpatrick, auditoría de
solapamiento con Derm1M mediante hash MD5).

\paragraph{O2. Construcción del dataset DermapixelAI 1.0.}
Se extrae, depura y estructura el material del archivo Dermapixel
para conformar un dataset que vincula, para cada caso, imagen y
modalidad, metadatos clínicos y texto narrativo en español. El
dataset se publica bajo licencia CC-BY-NC-SA~4.0 con datasheet
formal~\cite{datasheets2018}, particiones documentadas y mecanismo
de cita recomendada. Sobre él se ejecuta además parte de la
evaluación experimental del capítulo~\ref{cap:resultados}, que
contrasta el sondeo lineal supervisado frente a la clasificación
zero-shot multimodal en los tres niveles ontológicos
(secciones~\ref{sec:dermapixel-results} y~\ref{sec:dermapixel-zs}).

\paragraph{O3. Inicialización de Dermapixel R0, una primera línea
base dermatológica en español.}
A partir del codificador PanDerm Large y del dataset
DermapixelAI~1.0, se plantea adaptar un modelo fundacional
preexistente al castellano mediante adaptación supervisada de bajo
rango, dando lugar a una versión inicial (R0). Se describirá su
comportamiento sobre los niveles ontológicos L2 y L3
(sección~\ref{sec:dermapixel-spanderm-v0}) junto con una rama
contrastiva multimodal en español
(sección~\ref{sec:spanderm-clip}). El propósito no es alcanzar un
rendimiento de referencia, sino disponer de una semilla reproducible
y acotada que sirva de punto de partida para el escalado descrito en
la sección~\ref{sec:linea-spanderm}.

Los bloques derivados emergidos durante el desarrollo del proyecto
(armonización ontológica jerárquica L1/L2/L3, ensemble de seguridad
para detección de melanoma, análisis de equidad por fototipo,
interpretabilidad mediante Sparse Autoencoders~\cite{templeton2024},
recuperación visual densa basada en FAISS~\cite{faiss2019},
integración en el prototipo DermApIxel) no constituyen objetivos
formales y se documentan en los anexos y secciones correspondientes
con su estado, prerrequisitos y líneas futuras.


\section{Estructura del documento}
\label{sec:estructura}

El documento se compone de siete capítulos en el cuerpo principal,
una bibliografía y un anexo que describe el repositorio reproducible.
El capítulo~\ref{cap:sota} sintetiza
el estado del arte en cuatro bloques: aprendizaje profundo aplicado
a dermatología, modelos fundacionales en imagen médica, modelos
específicos para dermatología y brechas operativas no resueltas por
la literatura. El capítulo~\ref{cap:metodos} describe materiales y
métodos: los doce datasets externos, la ontología jerárquica
L1/L2/L3, el dataset DermapixelAI~1.0, los catorce modelos fundacionales
evaluados, la infraestructura, los protocolos, las métricas y la
trazabilidad. El
capítulo~\ref{cap:resultados} reúne los resultados experimentales
organizados por bloque. El capítulo~\ref{cap:discusion} interpreta
los resultados en clave de hallazgos, comparación con la literatura
e implicaciones operativas. El capítulo~\ref{cap:limitaciones}
declara las limitaciones del estudio. El
capítulo~\ref{cap:conclusiones} sintetiza el balance del trabajo y
enumera las contribuciones reproducibles. El material de soporte
---mapa de tareas experimentales, tablas detalladas de resultados por
dataset, pipeline de construcción del dataset DermapixelAI~1.0,
documento de entrega formal del dataset y ablaciones
complementarias--- se recoge en el repositorio reproducible descrito
en el anexo correspondiente.
